\section{A subthreshold minimizer on the cycle}
\label{sec:minimizer}

Throughout this section, $n>2$ and $A$ is the extremizing cycle constructed
in Lemma~\ref{lem:extremizer}, with the normalization
\eqref{eq:normalization}.  Thus $\mu=\|A\|$ is a finite measure on the
compact space $Y$, and
\[
 0<m:=\mu(Y)<s_n.
\]
We use the charts, multiplicities, and covector norms of
Proposition~\ref{prop:chart-mass}.  In particular, the chart images $Z_i$
are pairwise disjoint, their multiplicities
$\theta_i$ are nonzero signed integers, and
\[
 \mu=\sum_i(\phi_i)_\#\bigl(|\theta_i|J_i\,dz\bigr),
 \qquad J_i=\sqrt{\det G_i}.
\]
Set, as in Section~\ref{sec:sobolev},
\[
 \beta=\frac1{n-2},\qquad p=\frac{2n}{n-2},\qquad
 S_*=ns_n^{2/n},\qquad E(u)=\int_Y|Du|^2\,d\mu.
\]
All function norms in this section are taken with respect to $\mu$.
We use the almost sharp estimate \eqref{eq:almost-sobolev}; in particular,
its constant $C_\epsilon$ is independent of the Lipschitz function to
which it is applied.

We seek a nonnegative minimizer of
\[
 \frac{4\beta E(u)+n\|u\|_2^2}{\|u\|_p^2},
\]
initially defined for Lipschitz functions with $\|u\|_p>0$.
To pass to a minimizing limit, we first close the chart gradient and
prove compactness in $L^2$. The almost sharp inequality from
Section~\ref{sec:sobolev} then prevents loss of critical $L^p$ mass
below $S_*$.

\subsection{The closed gradient and calculus}

\begin{proposition}[The Sobolev space and its calculus]
\label{prop:sobolev-space}
Identify real Lipschitz functions on $Y$ that agree $\mu$-almost
everywhere. Their completion in the norm
\[
 \|u\|_{\Hcal}^2=\|u\|_2^2+E(u)
\]
is a Hilbert space $\Hcal$ of functions in $L^2(\mu)$.  Its gradient is
a well-defined closed operator with values in the Hilbert space of
square-integrable chart covectors.  The inclusion
$\Hcal\longrightarrow L^p(\mu)$ is continuous, and
\eqref{eq:almost-sobolev} holds for all $u\in\Hcal$.

The following calculus properties hold.
\begin{enumerate}[label=\textup{(\roman*)}]
 \item If $u_1,\ldots,u_k\in\Hcal$ and
 $\Phi\in C^1(\R^k)$ has bounded first derivatives, then
 $\Phi(u_1,\ldots,u_k)\in\Hcal$ and
 \[
 D\Phi(u_1,\ldots,u_k)
 =\sum_{a=1}^k\partial_a\Phi(u_1,\ldots,u_k)Du_a.
 \]
 \item For every $u\in\Hcal$ and every fixed $t\in\R$,
 $Du=0$ almost everywhere on $\{u=t\}$.  Consequently, the scalar
 chain rule also holds for Lipschitz, piecewise $C^1$ functions with
 finitely many corners, with any derivative values assigned at the
 corners.
 \item The map $u\mapsto u_+$ is continuous from $\Hcal$ to itself.
 Every nonnegative element of $\Hcal$ has nonnegative Lipschitz
 approximants converging in $\Hcal$.
 \item For each fixed center $y\in Y$, the radial inequality
 \eqref{eq:radial} holds for every nonnegative $g\in\Hcal$, with
 $r=d(y,\cdot)$ and its chart gradient $Dr$.
\end{enumerate}
\end{proposition}

\begin{proof}
Let $\mathcal L^2$ denote the Hilbert space of chart covector fields
$L$ with $\int|L|^2\,d\mu<\infty$. The Lipschitz level-set property
from Section~\ref{sec:currents}, applied to the difference of two
representatives, shows that their gradients agree whenever the functions
agree $\mu$-almost everywhere. We first prove that the graph of
the Lipschitz gradient in $L^2(\mu)\times\mathcal L^2$ is closable.
Suppose that $u_j$ are Lipschitz,
\[
 u_j\longrightarrow0\quad\hbox{in }L^2(\mu),\qquad
 Du_j\longrightarrow L\quad\hbox{in }\mathcal L^2.
\]
For a Lipschitz scalar function $b$ and a Lipschitz
$(n-1)$-tuple $\pi$, the identity $\partial A=0$ gives
\begin{equation}
 A(b,u_j,\pi)=-A(u_j,b,\pi).
 \label{eq:cycle-gradient-identity}
\end{equation}
The right-hand side tends to zero by the mass bound and
$\|u_j\|_1\le m^{1/2}\|u_j\|_2$.  For the fixed tuple $\pi$, the limit
of the left-hand side is integration of $b$ against the signed measure
\[
 \nu_\pi(B)=\sum_i
 \int_{\phi_i^{-1}(B\cap Z_i)}
 \theta_i\det\bigl(L_i,D(\pi\circ\phi_i)\bigr)\,dz,
 \qquad B\subset Y\ \hbox{Borel}.
\]
Here $L_i$ is the row of coordinate components of $L$ in the $i$th
chart.  The determinant estimate of
Proposition~\ref{prop:chart-mass} shows that this is a finite signed
measure and that
\[
 \|\nu_\pi\|_{\TV}
 \le \left(\prod_{a=1}^{n-1}\Lip(\pi_a)\right)
       \int|L|\,d\mu.
\]
The same estimate applied to $Du_j-L$ justifies the asserted limit.
Thus $\int b\,d\nu_\pi=0$ for every Lipschitz $b$ on $Y$.
Lipschitz functions are uniformly dense in $C(Y)$, so they determine
finite signed measures on the compact metric space $Y$.  It follows
that $\nu_\pi=0$ as a measure.

Fix a chart $i$.  Extend each coordinate of $\phi_i^{-1}:Z_i\to\R^n$
to a Lipschitz scalar function $F^a$ on $Y$.  At almost every density
point of $E_i$,
\[
 D(F^a\circ\phi_i)=e_a^*.
\]
Take for $\pi$ each of the $n$ tuples obtained by omitting one member
of $(F^1,\ldots,F^n)$.  Restricting the corresponding zero measures
$\nu_\pi$ to $Z_i$ shows, component by component, that
$\theta_i L_i=0$ almost everywhere on $E_i$.  The multiplicity is
nonzero, so $L_i=0$.  There are only countably many charts and finitely
many tuples per chart; hence $L=0$ in $\mathcal L^2$.
Notice the order of this argument: we first obtained zero signed
measures by testing on the whole cycle, and then localized those
measures.  No boundary or normality property of a chart restriction
is used.

The closure of the graph is therefore the graph of an operator $D$,
and, with its graph norm, it is the asserted Hilbert space $\Hcal$.
In particular, functions that agree almost everywhere have the same
gradient.  Applying \eqref{eq:almost-sobolev} to differences of
Lipschitz approximants shows that a graph-norm Cauchy sequence is
Cauchy in $L^p$.  Its $L^p$ limit agrees with its $L^2$ limit, and
passing to the limit proves both the continuous inclusion and the
extension of \eqref{eq:almost-sobolev}.

To prove the first calculus assertion, choose Lipschitz approximants
$u_{a,j}\to u_a$ in $\Hcal$, and pass to a common subsequence converging
almost everywhere.  Write $U_j=(u_{1,j},\ldots,u_{k,j})$ and
$U=(u_1,\ldots,u_k)$.  Bounded first derivatives give
$\Phi(U_j)\to\Phi(U)$ in $L^2$.  The constant $\Phi(0)$ is harmless
because $\mu$ is finite.  The chart chain rule gives
\[
 D\Phi(U_j)=\sum_a\partial_a\Phi(U_j)Du_{a,j}.
\]
For each summand, subtract the proposed limit by writing
\[
 \partial_a\Phi(U_j)(Du_{a,j}-Du_a)
 +\bigl(\partial_a\Phi(U_j)-\partial_a\Phi(U)\bigr)Du_a.
\]
The first term tends to zero in $\mathcal L^2$ by the derivative
bound; the second does so by dominated convergence.  Closedness
proves assertion~\textup{(i)}.

For assertion~\textup{(ii)}, fix a level $t$ and choose
$\eta\in C_c^\infty((-1,1))$ with $0\le\eta\le1$ and $\eta(0)=1$.
The functions
\[
 \Phi_\epsilon(s)=\int_{-\infty}^s
                  \eta\bigl((a-t)/\epsilon\bigr)\,da
\]
converge uniformly to zero, their derivatives are bounded by one,
and $\Phi_\epsilon'(s)\to\mathbf1_{\{t\}}(s)$.  Assertion~\textup{(i)}
and dominated convergence give
\[
 \Phi_\epsilon(u)\longrightarrow0\quad\hbox{in }L^2,
 \qquad
 D\Phi_\epsilon(u)\longrightarrow\mathbf1_{\{u=t\}}Du
 \quad\hbox{in }\mathcal L^2.
\]
Closedness forces the latter limit to be zero.  A Lipschitz,
piecewise $C^1$ scalar function with finitely many corners can now be
smoothed by convolution.  The smoothed functions converge uniformly,
their derivatives have a common bound, and the derivatives converge
at every point other than the corners.  The level-set conclusion
removes the contribution of those corners, so dominated convergence
and closedness give the claimed chain rule.

In particular, $D(u_+)=\mathbf1_{\{u>0\}}Du$.  If $u_j\to u$ in
$\Hcal$, the positive parts converge in $L^2$, and their gradient
difference is
\[
 \mathbf1_{\{u_j>0\}}(Du_j-Du)
 +\bigl(\mathbf1_{\{u_j>0\}}-\mathbf1_{\{u>0\}}\bigr)Du.
\]
Along any subsequence with almost-everywhere convergence, the second
term tends to zero in $\mathcal L^2$ by the level-set conclusion at
zero and dominated convergence.  The first term is bounded in norm
by $\|Du_j-Du\|_2$.  Applying this observation to subsequences proves
continuity for the original sequence.  Taking positive parts of
Lipschitz approximants proves assertion~\textup{(iii)}.

Finally fix $y\in Y$.  The distance $r=d(y,\cdot)$ is bounded,
$|Dr|\le1$ almost everywhere, and $\mu$ is finite.  Each term of
\eqref{eq:radial} is therefore continuous as a linear functional of
$g$ in the graph norm.  Approximate a nonnegative $g\in\Hcal$ by the
nonnegative Lipschitz functions just constructed and pass to the
limit.  This proves assertion~\textup{(iv)}.
\end{proof}

\subsection{Compactness}

The gradient is now defined on the completed space, and its critical
embedding is continuous. We still need strong $L^2$ convergence for
bounded sequences. The proof obtains compactness for projections of
whole weighted currents, then isolates charts in total variation.

\begin{proposition}[Compact inclusion]
\label{prop:compact-embedding}
The inclusion $\Hcal\longrightarrow L^2(\mu)$ is compact.
\end{proposition}

\begin{proof}
\emph{Projecting the whole weighted current.}
First consider Lipschitz functions $u_j$ with
$\sup_j\|u_j\|_{\Hcal}\le M_0<\infty$.  Fix a chart $i$, extend its
inverse coordinates to a Lipschitz map $F=(F^1,\ldots,F^n):Y\to\R^n$,
and fix a Lipschitz function $\chi$ on $Y$.  Define finite signed
measures $\lambda_{j,\chi}$ on $\R^n$ by
\begin{equation}
 \int h\,d\lambda_{j,\chi}
 =A\bigl((h\circ F)\chi u_j,F^1,\ldots,F^n\bigr)
 \label{eq:projected-function-measure}
\end{equation}
for bounded Borel $h$.  The chart representation defines these
measures directly.  They are supported in the fixed compact set
$F(Y)$, and the current mass bound gives
\[
 \sup_j\|\lambda_{j,\chi}\|_{\TV}
 \le \left(\prod_{a=1}^n\Lip(F^a)\right)
       \|\chi\|_\infty m^{1/2}M_0.
\]

For $h\in C_c^\infty(\R^n)$, expanding the row
$D(h\circ F)=\sum_a(\partial_a h\circ F)DF^a$ in a determinant shows
that
\[
 \int\partial_l h\,d\lambda_{j,\chi}
 =A\bigl(\chi u_j,F^1,\ldots,F^{l-1},h\circ F,
                         F^{l+1},\ldots,F^n\bigr).
\]
All terms except $a=l$ in the row expansion vanish because they have
a repeated row.  Apply the cycle identity
\eqref{eq:cycle-gradient-identity}, with a permutation of rows if
necessary, to move $h\circ F$ to the first coefficient.  The intrinsic
determinant bound then gives
\begin{equation}
 \left|\int\partial_lh\,d\lambda_{j,\chi}\right|
 \le C_F\|h\|_\infty\int|D(\chi u_j)|\,d\mu
 \le C_\chi\|h\|_\infty,
 \label{eq:projected-measure-derivative}
\end{equation}
uniformly in $j$ and $l$.  Indeed,
\[
 \int|D(\chi u_j)|\,d\mu
 \le m^{1/2}\bigl(\|\chi\|_\infty\|Du_j\|_2
                      +\Lip(\chi)\|u_j\|_2\bigr).
\]
Constants here may depend on the fixed chart, $\chi$, and $M_0$.

\emph{Translation compactness.}
We give the compactness consequence of
\eqref{eq:projected-measure-derivative} in detail. This is the usual
translation and mollification proof of Euclidean BV compactness
\cite[Chapter~2, Theorem~2.6]{Simon2014}, expressed directly for the
projected signed measures.  Let
$\tau_w(x)=x+w$.  For a smooth compactly supported test $h$, the
fundamental theorem of calculus along $x+tw$ and
\eqref{eq:projected-measure-derivative} imply
\[
 \left|\int h\,d\bigl((\tau_w)_\#\lambda_{j,\chi}
                              -\lambda_{j,\chi}\bigr)\right|
 \le C_\chi\|h\|_\infty\sum_{l=1}^n|w_l|.
\]
Smooth compactly supported tests recover the total variation of
finite signed Radon measures, by regularity and approximation of
continuous tests.  After changing the constant, we obtain
\begin{equation}
 \| (\tau_w)_\#\lambda_{j,\chi}-\lambda_{j,\chi}\|_{\TV}
 \le C_\chi|w|.
 \label{eq:measure-translation-bound}
\end{equation}
Choose a nonnegative smooth convolution kernel $\rho_\epsilon$ of
integral one supported in the ball of radius $\epsilon$.  Averaging
\eqref{eq:measure-translation-bound} gives
\[
 \|\lambda_{j,\chi}-(\rho_\epsilon*\lambda_{j,\chi})\,dx\|_{\TV}
 \le C_\chi\epsilon.
\]
For fixed $\epsilon$, the densities
$\rho_\epsilon*\lambda_{j,\chi}$ have a common compact support and
uniform bounds on their values and first derivatives.  The
Arzel\`a--Ascoli theorem makes them precompact in the uniform norm
on that support and hence in $L^1(\R^n)$.  Approximation by these
families, first fixing $\epsilon$ and then letting it decrease to
zero, proves that $\{\lambda_{j,\chi}:j\ge1\}$ is totally bounded
in total variation.  The space of finite signed measures is complete
in that norm, so this family is precompact.

\emph{Isolating a chart in total variation.}
We now isolate the chart.  Formula
\eqref{eq:projected-function-measure} also defines
$\lambda_{j,\mathbf1_{Z_i}}$, using the bounded Borel coefficient
$\mathbf1_{Z_i}$.  By regularity of $\mu$ and distance cutoffs on the
compact metric space $Y$, choose Lipschitz $\chi_k$ with
\[
 \|\chi_k-\mathbf1_{Z_i}\|_2\longrightarrow0.
\]
The mass bound and Cauchy--Schwarz give, uniformly in $j$,
\begin{align}
 \|\lambda_{j,\chi_k}-\lambda_{j,\mathbf1_{Z_i}}\|_{\TV}
 &\le \left(\prod_{a=1}^n\Lip(F^a)\right)
       \int|\chi_k-\mathbf1_{Z_i}|\,|u_j|\,d\mu \notag\\
 &\le C_FM_0\|\chi_k-\mathbf1_{Z_i}\|_2.
 \label{eq:chart-cutoff-approximation}
\end{align}
For each fixed $k$ the family on the left with coefficient $\chi_k$
is precompact by the preceding argument.  Uniform approximation
\eqref{eq:chart-cutoff-approximation} therefore makes
$\{\lambda_{j,\mathbf1_{Z_i}}:j\ge1\}$ precompact as well.  There is
no need for the Lipschitz constants of $\chi_k$ to remain bounded.

Since $F\circ\phi_i$ is the identity on $E_i$, the density of the
isolated measure is exactly
\[
 d\lambda_{j,\mathbf1_{Z_i}}(z)
   =\mathbf1_{E_i}(z)\theta_i(z)(u_j\circ\phi_i)(z)\,dz.
\]
Consequently, for any $j,k$,
\begin{align*}
 \|u_j-u_k\|_{L^1(Z_i,\mu)}
 &=\int_{E_i}|\theta_i|J_i
                  |(u_j-u_k)\circ\phi_i|\,dz\\
 &\le\Lip(\phi_i)^n
          \|\lambda_{j,\mathbf1_{Z_i}}
                    -\lambda_{k,\mathbf1_{Z_i}}\|_{\TV}.
\end{align*}
Here $J_i\le\Lip(\phi_i)^n$.  This calculation allows arbitrary
signed, unbounded integer multiplicities: the fixed $\theta_i$
appears as $|\theta_i|$ in the total variation of a difference.
Likewise, \eqref{eq:chart-cutoff-approximation} already uses the full
mass measure $|\theta_i|J_i\,dz$, so it requires no bound or regularity
of the multiplicities.  All integrations by parts above concerned
the whole cycle $A$ with a Lipschitz coefficient.  None concerned a
restriction to one chart.

\emph{Assembling the charts.}
We may now take a diagonal subsequence that is Cauchy in
$L^1(Z_i,\mu)$ for every $i$.  The omitted charts have uniformly
small contributions, because
\[
 \int_{\bigcup_{i>N}Z_i}|u_j|\,d\mu
 \le M_0\,\mu\left(\bigcup_{i>N}Z_i\right)^{1/2}
 \longrightarrow0
\]
uniformly in $j$.  The chart images cover $\mu$ up to a null set,
and $\mu$ is finite.  The subsequence is thus Cauchy in $L^1(\mu)$.
Proposition~\ref{prop:sobolev-space} supplies a uniform $L^p$ bound.
Since $p>2$, interpolation with
$\alpha=(p-2)/(2(p-1))>0$ gives
\[
 \|u_j-u_k\|_2
 \le\|u_j-u_k\|_1^\alpha\|u_j-u_k\|_p^{1-\alpha}
 \longrightarrow0.
\]
Finally, approximate the $j$th member of an arbitrary bounded
sequence in $\Hcal$ by a Lipschitz function with graph-norm error
at most $1/j$.  The result just proved for those approximants proves
compactness for the original sequence.
\end{proof}

\subsection{A minimizer and its powers}

We now have weak compactness in $\Hcal$ and strong compactness in
$L^2$. These do not alone give strong convergence at the critical
exponent $p$. The strict inequality $m<s_n$ places the quotient
infimum below $S_*$; the almost sharp estimate will exclude any missing
$L^p$ mass.

Define the quadratic functional and its critical constrained infimum by
\begin{equation}
 Q(u)=4\beta E(u)+n\|u\|_2^2,
 \qquad
 \Lambda=\inf\{Q(u):u\in\Hcal,\ \|u\|_p=1\}.
 \label{eq:critical-variational-problem}
\end{equation}

\begin{proposition}[Existence and normalization of the minimizer]
\label{prop:minimizer}
One has $0<\Lambda<S_*$.  There exists a nonzero nonnegative
$v\in\Hcal$ minimizing $Q(u)/\|u\|_p^2$ over nonzero $u\in\Hcal$
and satisfying
\begin{equation}
 \begin{split}
 0<W:=\int v^p\,d\mu
       =\left(\frac\Lambda n\right)^{n/2}<s_n,
 \qquad Q(v)=nW.
 \end{split}
 \label{eq:minimizer-normalization}
\end{equation}
For every $\psi\in\Hcal$ it satisfies the weak equation
\begin{equation}
 4\beta\int Dv\cdot D\psi\,d\mu+n\int v\psi\,d\mu
   =n\int v^{p-1}\psi\,d\mu.
 \label{eq:euler}
\end{equation}
\end{proposition}

\begin{proof}
Taking $\epsilon=S_*/2$ in \eqref{eq:almost-sobolev} shows that
\[
 \frac{S_*}{2}\|u\|_p^2
 \le4\beta E(u)+C_{S_*/2}\|u\|_2^2
 \le\max\{1,C_{S_*/2}/n\}\,Q(u).
\]
Thus $\Lambda>0$.  The constant function $m^{-1/p}$ has $L^p$ norm
one, so
\begin{equation}
 \Lambda\le n m^{1-2/p}=n m^{2/n}<ns_n^{2/n}=S_*.
 \label{eq:strict-subthreshold}
\end{equation}

Let $u_j$ be an $L^p$-normalized minimizing sequence.  It is bounded
in $\Hcal$, since $Q$ is a positive definite quadratic form equivalent
to the squared graph norm.  By weak compactness in a Hilbert space
and Proposition~\ref{prop:compact-embedding}, after taking a
subsequence we have
\[
 u_j\rightharpoonup u\quad\hbox{in }\Hcal,
 \qquad u_j\to u\quad\hbox{in }L^2(\mu)
 \quad\hbox{and almost everywhere}.
\]
Put $z_j=u_j-u$ and $t=\int|u|^p\,d\mu$.  Fatou's lemma gives
$0\le t\le1$. We give the Br\'ezis--Lieb
splitting~\cite[Theorem~1]{BrezisLieb1983} in the form needed here:
\begin{equation}
 \int|z_j|^p\,d\mu\longrightarrow1-t.
 \label{eq:critical-mass-splitting}
\end{equation}
For every $\epsilon>0$, the mean-value theorem and Young's
inequality give a finite $C_\epsilon$ such that, for real $a,b$,
\[
 \bigl||a+b|^p-|a|^p\bigr|
 \le\epsilon|a|^p+C_\epsilon|b|^p.
\]
Apply this with $a=z_j$ and $b=u$, and set
$R_j=|u_j|^p-|z_j|^p-|u|^p$.  The functions
\[
 \bigl(|R_j|-\epsilon|z_j|^p\bigr)_+
\]
converge to zero almost everywhere and are bounded by
$(C_\epsilon+1)|u|^p$.  Dominated convergence, followed by
$\epsilon\downarrow0$, shows that $R_j\to0$ in $L^1$, since the
$L^p$ norms of $z_j$ are bounded.  This proves
\eqref{eq:critical-mass-splitting}.

Weak convergence in $\Hcal$ also gives the quadratic splitting
\[
 Q(u_j)=Q(u)+Q(z_j)+o(1).
\]
The definition of $\Lambda$ gives $Q(u)\ge\Lambda t^{2/p}$, with
the same conclusion if $u=0$.  For fixed $\epsilon>0$, apply
\eqref{eq:almost-sobolev} to $z_j$ and use $\|z_j\|_2\to0$ to obtain
\[
 \liminf_j Q(z_j)\ge(S_*-\epsilon)(1-t)^{2/p}.
\]
Letting $\epsilon\downarrow0$ therefore yields
\begin{equation}
 \Lambda\ge\Lambda t^{2/p}+S_*(1-t)^{2/p}.
 \label{eq:no-critical-mass-loss}
\end{equation}
If $t<1$, the inequalities $a^{2/p}\ge a$ for $0\le a\le1$ and
$S_*>\Lambda$ make the right-hand side strictly larger than
$\Lambda$: indeed it is at least
$\Lambda t+S_*(1-t)>\Lambda$.  Thus $t=1$.
Weak lower semicontinuity now gives $Q(u)\le\Lambda$, so $u$ is a
minimizer with $\|u\|_p=1$.  Replacing it by $|u|$ preserves both
its $L^p$ norm and its energy, by
Proposition~\ref{prop:sobolev-space}.  We take $u\ge0$.

For each $\psi\in\Hcal$, differentiate the quotient
$Q(u+s\psi)/\|u+s\psi\|_p^2$ at $s=0$.  Its denominator is nonzero
for sufficiently small $s$, and H\"older's inequality justifies the
derivative of the $L^p$ integral: for $|s|\le1$ the derivative
integrand is bounded by
$p(|u|+|\psi|)^{p-1}|\psi|\in L^1$.  Since $u$ is a minimizer and
$\|u\|_p=1$, this gives
\[
 4\beta\int Du\cdot D\psi\,d\mu+n\int u\psi\,d\mu
 =\Lambda\int u^{p-1}\psi\,d\mu.
\]
Set
\[
 v=\left(\frac\Lambda n\right)^{1/(p-2)}u.
\]
Homogeneity preserves the quotient minimum, and
$\Lambda(\Lambda/n)^{-1}=n$ gives \eqref{eq:euler}.
Since $p/(p-2)=n/2$, its $p$th moment is
$W=(\Lambda/n)^{n/2}<s_n$ by
\eqref{eq:strict-subthreshold}.  Testing \eqref{eq:euler} with
$\psi=v$ gives $Q(v)=nW$, completing
\eqref{eq:minimizer-normalization}.
\end{proof}

The normalized minimizer has the strict weighted mass bound $W<s_n$
needed for the contradiction. The remaining task is to justify nonlinear
tests in its weak equation and radial first variation. The following
moment and weighted-energy estimates do so without assuming that $v$
is positive everywhere, smooth, or essentially bounded.
The proof adapts the critical-potential absorption method of
Br\'ezis and Kato~\cite{BrezisKato1979} to the closed chart gradient.
For the truncated-power formulation, see also \cite[Lemma~6, p.~178]{Brezis1986}.
The adaptation is given in full, without invoking a Euclidean regularity
theorem in the present Sobolev space.

\begin{lemma}[Powers of the minimizer]
\label{lem:powers}
For every real $\gamma\ge1$,
\[
 v^\gamma\in\Hcal\cap L^p(\mu),\qquad
 D(v^\gamma)=\gamma v^{\gamma-1}Dv.
\]
Moreover, $\int v^q\,d\mu<\infty$ for every finite $q>0$, and
\begin{equation}
 \int v^q|Dv|^2\,d\mu<\infty
 \qquad\hbox{for every finite }q\ge0.
 \label{eq:polynomial-weighted-energy}
\end{equation}
\end{lemma}

\begin{proof}
We first prove the power assertion under the additional assumption
$v\in L^{2\gamma}$.  It is immediate for $\gamma=1$, so suppose
$\gamma>1$.  For $N\ge1$ define
\[
 w_N=v\min(v,N)^{\gamma-1},\qquad
 \psi_N=v\min(v,N)^{2\gamma-2},\qquad
 d_\gamma=\frac{2\gamma-1}{\gamma^2}.
\]
The scalar functions in these definitions, extended by zero for
negative arguments, are Lipschitz and piecewise $C^1$.  Thus both
compositions belong to $\Hcal$.  Their derivatives show that
\[
 Dv\cdot D\psi_N\ge d_\gamma|Dw_N|^2.
\]
Indeed there is equality below $N$, while above $N$ the ratio of
the coefficients is one and $0<d_\gamma\le1$.  The gradient vanishes
on the level set $\{v=N\}$, so no value at the corner matters.
Testing \eqref{eq:euler} with $\psi_N$ is legitimate; for example,
$v^{p-1}\psi_N\le N^{2\gamma-2}v^p$ is integrable.  Dropping the
nonnegative term $n\int v\psi_N\,d\mu$ gives
\begin{equation}
 4\beta d_\gamma E(w_N)
 \le n\int v^{p-2}w_N^2\,d\mu.
 \label{eq:truncated-power-energy}
\end{equation}

Apply \eqref{eq:almost-sobolev} to $w_N$ with $\epsilon=S_*/2$ and
multiply by $d_\gamma$.  By
\eqref{eq:truncated-power-energy},
\[
 \frac{d_\gamma S_*}{2}\|w_N\|_p^2
 \le n\int v^{p-2}w_N^2\,d\mu
       +d_\gamma C_{S_*/2}\|w_N\|_2^2.
\]
H\"older's inequality with conjugate exponents $p/(p-2)$ and $p/2$
gives
\[
 n\int_{\{v>K\}}v^{p-2}w_N^2\,d\mu
 \le n\left(\int_{\{v>K\}}v^p\,d\mu\right)^{(p-2)/p}
       \|w_N\|_p^2.
\]
Since $v\in L^p$, choose $K\ge1$, depending on $\gamma$ but not $N$,
so that the coefficient on the right is at most
$d_\gamma S_*/4$.  The contribution of $\{v\le K\}$ is at most
$nK^{p-2}\|w_N\|_2^2$.  Absorbing the tail yields
\[
 \frac{d_\gamma S_*}{4}\|w_N\|_p^2
 \le\bigl(nK^{p-2}+d_\gamma C_{S_*/2}\bigr)\|w_N\|_2^2.
\]
The assumed $L^{2\gamma}$ moment bounds the right-hand side uniformly
in $N$, because $w_N\le v^\gamma$.  Thus $\|w_N\|_p$ is uniformly
bounded.  Equation~\eqref{eq:truncated-power-energy}, with another
application of H\"older, then uniformly bounds $E(w_N)$.

On $\{v<N\}$ the gradient of $w_N$ is
$\gamma v^{\gamma-1}Dv$.  Monotone convergence in these sets shows
that $\gamma v^{\gamma-1}Dv$ belongs to $\mathcal L^2$.  At every
truncation level the formulas and the level-set property give
\[
 |w_N|\le v^\gamma,\qquad
 |Dw_N|\le\gamma v^{\gamma-1}|Dv|.
\]
The values and gradients converge almost everywhere to $v^\gamma$
and $\gamma v^{\gamma-1}Dv$, respectively.  The displayed bounds
give strong $L^2$ convergence of both by dominated convergence.
Closedness proves $v^\gamma\in\Hcal$ with the asserted gradient,
and the continuous inclusion into $L^p$ gives $v^\gamma\in L^p$.

Starting with $v\in L^p$, set $\ell_j=p(p/2)^j$ for $j\ge0$.
The conditional assertion just proved, with $\gamma=\ell_j/2$, shows
\[
 v\in L^{\ell_j}\quad\Longrightarrow\quad v\in L^{\ell_{j+1}}.
\]
Since $p/2>1$, these exponents tend to infinity.  Finiteness of $\mu$ then gives
every finite positive moment of $v$.  The conditional assertion
therefore applies to every $\gamma\ge1$.  Finally, for $q\ge0$
take $\gamma=1+q/2$ to obtain
\[
 \int v^q|Dv|^2\,d\mu=\gamma^{-2}E(v^\gamma)<\infty,
\]
which is \eqref{eq:polynomial-weighted-energy}.
\end{proof}

All positive moments and polynomially weighted energies are now
finite. We next turn those integrability statements into an admissible
chain rule for the actual two-variable tests used in the chord
comparison.

\begin{lemma}[Compositions with polynomial growth]
\label{lem:composites}
Let $r\in\Hcal$ satisfy $0\le r\le R<\infty$ almost everywhere and
$|Dr|\in L^\infty(\mu)$.  Suppose that $F$ is $C^1$ on an open
neighborhood of $[0,R]\times[0,\infty)$ and that, on this strip,
\[
 |F(a,t)|+|\partial_1F(a,t)|+|\partial_2F(a,t)|
 \le C(1+t)^M
\]
for some finite $C$ and nonnegative integer $M$.  Then
$F(r,v)\in\Hcal$ and
\begin{equation}
 D(F(r,v))=\partial_1F(r,v)Dr+\partial_2F(r,v)Dv.
 \label{eq:polynomial-composite-chain-rule}
\end{equation}
In particular, this applies with $r=d(y,\cdot)$ for each fixed
$y\in Y$.
\end{lemma}

\begin{proof}
Let $t_N=\min(v,N)$.  The scalar calculus in
Proposition~\ref{prop:sobolev-space} gives $t_N\in\Hcal$ and
$Dt_N=\mathbf1_{\{v<N\}}Dv$, where the level-set property handles
$v=N$.  For each $N$, choose a smooth, compactly supported cutoff in the given
open neighborhood and equal to one on a neighborhood of the compact
rectangle $[0,R]\times[0,N]$.  Multiply $F$ by this cutoff and extend
by zero to obtain a global $C^1$ function with bounded first
derivatives.  It agrees with $F$ and its first derivatives on the
rectangle.  The bounded-derivative calculus consequently gives
$F(r,t_N)\in\Hcal$ and
\[
 D(F(r,t_N))
 =\partial_1F(r,t_N)Dr
  +\partial_2F(r,t_N)\mathbf1_{\{v<N\}}Dv.
\]
Writing $L=\|\,|Dr|\,\|_\infty$, the growth assumption gives bounds
independent of $N$,
\[
 |F(r,t_N)|\le C(1+v)^M,\qquad
 |D(F(r,t_N))|\le C(1+v)^M(L+|Dv|).
\]
The squares of these bounds are integrable by
Lemma~\ref{lem:powers}: they are bounded by a constant times
$(1+v^{2M})(1+|Dv|^2)$.  Since $v$ is finite almost everywhere,
the values and the displayed gradient formulas converge almost
everywhere to the values and right-hand side in
\eqref{eq:polynomial-composite-chain-rule}.  Dominated convergence
gives strong convergence in the two $L^2$ spaces, and closedness
proves the result.  A distance function has bounded values and
$|Dr|\le1$, so it satisfies the hypotheses.
\end{proof}
